%% kapa: exam=1 round=23 year=2012 period=2 form=A subject=appraisal_law subject_ko=부동산관계법규 seq=16 number=56
%% source: Q-Net 감정평가사 기출문제 (cst00309) artlSeq=5052580 file=raw/1008/5052580/2012년 제23회 감정평가사 제1차시험 기출문제(2교시 A형).pdf
%% conversion: kordoc 4.24.51 -> tools/gen1.py (auto); review=proofread-sonnet; flags=none
\begin{question}{56}
국토의 계획 및 이용에 관한 법령상 용도지역에 관한 설명으로 옳지 않은 것은?
\begin{choices}
  \item 도시지역, 관리지역, 농림지역 또는 자연환경보전지역으로 용도가 지정되지 아니한 지역에서 용도지역에서의 행위제한 등에 대하여는 자연환경보전지역에 관한 규정을 적용한다.
  \item 용도지역이 미세분된 도시지역에서의 행위제한 등에 대하여는 보전녹지지역에 관한 규정을 적용한다.
  \item 관리지역에서「농지법」에 따른 농업진흥지역으로 지정ㆍ고시된 지역은 생산관리지역으로 결정ㆍ고시된 것으로 본다.
  \item 용도지역이 미세분된 관리지역에서의 행위제한 등에 대하여는 보전관리지역에 관한 규정을 적용한다.
  \item 바다인 공유수면의 매립목적이 그 매립구역과 이웃하고 있는 용도지역의 내용과 같으면, 도시관리계획의 입안 및 결정절차 없이 그 매립준공구역은 그 매립의 준공인가일부터 이와 이웃하고 있는 용도지역으로 지정된 것으로 본다.
\end{choices}
\end{question}
\answer{3}
